% from ega4-18-fr.tex
\oldpage[IV-4]{126}

Notons encore que, conformément à la théorie générale ($0_{\mathrm{III}}$, 8.1.6), l'automorphisme identique de $\mathbf V(\check{\sh{E}})$ correspond canoniquement à une section $c$ de $\sh{E}_{(\mathbf V(\check{\sh{E}}))}$ au-dessus de $\mathbf V(\check{\sh{E}})$, c'est-à-dire (II, 1.4.1) à un homomorphisme de $\mathbf S_{\sh{O}_X}(\check{\sh{E}})$-Modules
\[
  u:\mathbf S_{\sh{O}_X}(\check{\sh{E}})\longrightarrow
  \mathbf S_{\sh{O}_X}(\check{\sh{E}})\otimes_{\sh{O}_X}\sh{E};
\]
pour tout ouvert affine $W$ de $X$, si l'on pose $\Gamma(W,\sh{O}_X)=A$, si on identifie $\Gamma(W,\mathbf S_{\sh{O}_X}(\check{\sh{E}}))$ à une algèbre de polynômes $C=A[T_1,\ldots,T_n]$, de sorte que les $T_i$ forment une base de $\Gamma(W,\check{\sh{E}})$, et si enfin on désigne par $(e_i)$ la base duale de $(T_i)$ dans $\Gamma(W,\sh{E})=\Gamma(W,\check{\sh{E}})^\vee$, on voit aussitôt que $u$ correspond à l'homomorphisme de $C$-modules tel que
\[
  u(1)=\sum_{i=1}^n T_i\otimes e_i.
\]
Si $X'=\Spec(A)$ est affine et tel que $\sh{E}_{(X')}$ soit isomorphe à $\sh{O}_{X'}^n$, $\Gamma(X',\sh{E}_{(X')})$ est un $A$-module libre de rang $n$; de façon imagée, on peut dire que l'objet $\mathbf V(\check{\sh{E}})$ représente « l'ensemble des points de l'espace affine tordu sur $X$ défini par $\sh{E}$ ».

\begin{env}[18.5.2]
\label{IV.18.5.2}
Rappelons (II, 1.7.8) que l'on a par définition $\mathbf V(\check{\sh{E}})=\Spec(\mathbf S_{\sh{O}_X}(\check{\sh{E}}))$. Soit $\mathscr J$ un Idéal quasi-cohérent de $\mathbf S_{\sh{O}_X}(\check{\sh{E}})$, de sorte que $\Spec(\mathbf S_{\sh{O}_X}(\check{\sh{E}})/\mathscr J)$ est un sous-préschéma fermé de $\mathbf V(\check{\sh{E}})$; nous allons l'interpréter comme \emph{représentant un foncteur} de la catégorie des $X$-préschémas dans celle des ensembles. Notons pour cela qu'une section $u\in\Gamma(X,\sh{E})$ s'identifie canoniquement à un $\sh{O}_X$-homomorphisme $u:\sh{O}_X\to\sh{E}$, auquel correspond par transposition un $\sh{O}_X$-homomorphisme ${}^t u:\check{\sh{E}}\to\sh{O}_X$, et par suite un homomorphisme de $\sh{O}_X$-Algèbres $v:\mathbf S_{\sh{O}_X}(\check{\sh{E}})\to\sh{O}_X$. Soit $\operatorname{Al}(X,\sh{E},\mathscr J)$ l'ensemble des $u\in\Gamma(X,\sh{E})$ tels que $\mathscr J$ soit \emph{contenu dans le noyau} de $v$; il résulte aussitôt de ces définitions que
\[
  X'\longmapsto\operatorname{Al}(X',\sh{E}_{(X')},\mathscr J\otimes_{\sh{O}_X}\sh{O}_{X'})
\]
est un foncteur représenté par $\Spec(\mathbf S_{\sh{O}_X}(\check{\sh{E}})/\mathscr J)$. Si $X'=\Spec(A)$ est affine et tel que $\sh{E}_{(X')}$ soit isomorphe à $\sh{O}_{X'}^n$, $\Gamma(X',\sh{E}_{(X')})$ peut s'identifier à l'ensemble $A^n$, et $\mathscr J\otimes_{\sh{O}_X}\sh{O}_{X'}$ à un $\sh{O}_{X'}$-Module de la forme $\widetilde{\mathfrak J}$, où $\mathfrak J$ est un idéal de l'anneau de polynômes $A[T_1,\ldots,T_n]$; l'ensemble $\operatorname{Al}(X',\sh{E}_{(X')},\mathscr J\otimes_{\sh{O}_X}\sh{O}_{X'})$ s'identifie alors à la partie de $A^n$ formée des points $(t_1,\ldots,t_n)$ tels que $P(t_1,\ldots,t_n)=0$ pour tous les polynômes $P\in\mathfrak J$; de façon imagée on peut donc dire que l'objet $\Spec(\mathbf S_{\sh{O}_X}(\check{\sh{E}})/\mathscr J)$ représente « le sous-ensemble algébrique de l'espace affine tordu $\Gamma(X,\sh{E})$ formé des points annulant l'idéal $\Gamma(X,\mathscr J)$ ». On note encore ce $X$-préschéma $\mathbf{Al}(\sh{E},\mathscr J)$. On notera que si $\mathscr J$ est un Idéal de type fini de $\mathbf S_{\sh{O}_X}(\check{\sh{E}})$, $\mathbf{Al}(\sh{E},\mathscr J)$ est un $X$-préschéma \emph{de présentation finie}, puisque $\mathbf S_{\sh{O}_X}(\check{\sh{E}})$ est une $\sh{O}_X$-Algèbre de présentation finie.
\end{env}

\begin{lemma}[18.5.3]
\label{IV.18.5.3}
Soient $S$ un préschéma, $f:X\to S$ un morphisme fini et localement libre (18.2.7). Considérons le foncteur contravariant de la catégorie des $S$-préschémas dans la catégorie des ensembles
\begin{equation}
  S'\longmapsto\operatorname{Of}(X\times_S S')
  \tag{18.5.3.1}\label{IV.18.5.3.1}
\end{equation}
où $\operatorname{Of}(X\times_S S')$ est l'ensemble des parties à la fois ouvertes et fermées de l'espace sous-jacent à $X\times_S S'$. Alors ce foncteur est représentable par un $S$-préschéma $\mathbf{Of}(X)$, qui est affine, étale et de présentation finie sur $S$.
\end{lemma}

